\section{Introduction}
\label{sec:intro}

Over the past decade, emerging economies across South Asia and the Global South have experienced rapid growth in mobile application usage, establishing these platforms as critical infrastructure for financial transactions, welfare disbursements, and digital governance~\cite{chen_fintech_asia, worldbank_fintech_bd}. In jurisdictions such as Bangladesh, where active mobile financial wallets process billions of dollars annually, mobile endpoints represent a primary line of national cybersecurity and personal data defense~\cite{reaves_mfs}. In response to escalating cyber threats and unauthorized profiling, sovereign legislatures and regulators have promulgated statutory and regulatory frameworks, including the cabinet-approved \textit{Draft Personal Data Protection Act (PDPA)}~\cite{bd_pdpa_2026}, the \textit{Bangladesh Bank Cybersecurity Framework v1.0 (BB-CSF)}~\cite{bb_csf_2025}, and the \textit{Cyber Security Act, 2023 (CSA, Act No. XXVII of 2023)}~\cite{bd_csa_2026}.

These statutes mandate technical safeguards aligned with international verification standards such as OWASP MASVS~\cite{owasp_masvs}: transport-layer encryption, hardware-backed keystore segregation, cryptographic data destruction, least-privilege permission architectures, and explicit consent handling for personally identifiable information (PII). Non-conformance with enacted instruments exposes providers to administrative sanctions and license revocations, and the Draft PDPA, once enacted, will impose additional data-protection penalties. Consequently, software engineering teams, regulatory compliance officers, and independent auditors require automated verification methodologies to evaluate production mobile software against statutory mandates.

\subsection{The Epistemic Overreach Trap in Current Tools}
\label{sec:intro:dilemma}

Despite the need for regulatory verification, the software engineering community lacks formal frameworks to bridge natural-language legal requirements and low-level program behavior. In practice, auditors frequently repurpose generic mobile security scanners such as MobSF~\cite{mobsf} or Androguard~\cite{androguard}. While valuable for discovering generic vulnerabilities, repurposing these tools for statutory verification introduces four fundamental flaws:

\paragraph{1) The epistemic overreach trap} Conventional scanners routinely report scalar compliance percentages (e.g., ``78\% Compliant''). Such assertions commit a severe epistemic error~\cite{fagin_epistemic_logic}. An Android application package (APK/AAB) contains compiled Dalvik Executable (\texttt{.dex}) bytecode, an \texttt{AndroidManifest.xml} manifest, and packaged assets~\cite{enck_android_study}. While revealing client-side implementation details, an APK is an inherently partial observation of an enterprise system. An analyzer inspecting client bytecode cannot observe whether cloud databases are encrypted at rest, whether corporate officers respect retention schedules, or whether valid consent records are archived off-device. Conflating the absence of client-side findings with comprehensive organizational compliance constitutes epistemic hallucination.

\paragraph{2) Absence of formal traceability and provenance} Existing scanners emit disconnected vulnerability lists without a verifiable provenance chain linking low-level bytecode instructions to authoritative statutory clauses. When regulatory bodies or judicial authorities demand evidence, heuristic scanner outputs provide insufficient legal grounding.

\paragraph{3) Coupled and brittle rule encodings} Conventional scanners hardcode inspection logic directly inside analyzer scripts. When a regulatory body amends a legal clause, the underlying static analyzer must be refactored, violating separation-of-concerns principles.

\paragraph{4) Linguistic and jurisdictional alienation} Compliance outputs in existing tools are generated exclusively in English-language technical jargon. In jurisdictions where statutes are drafted in national languages (such as Bengali in Bangladesh) and software teams operate bilingually, this language barrier introduces cognitive friction, impairing defect interpretation and delaying remediation.

\subsection{The Reg2App Paradigm}
\label{sec:intro:paradigm}

To resolve these limitations, we propose \textbf{Reg2App}, a formally specified, provenance-preserving Regulation-as-Code methodology for translating statutory requirements into partially observable software properties. Reg2App is grounded in the principle of \textit{epistemic restraint}: an automated verification tool must explicitly delineate program properties provably observable from client artifacts from obligations residing beyond client-side observational boundaries~\cite{octeau_reflection}. Rather than synthesizing arbitrary compliance percentages, Reg2App formalizes an end-to-end translation pipeline:
\begin{align*}
\text{Statute} \rightarrow \text{Legal Obligation} \rightarrow \text{Technical Control} \rightarrow \text{Program Property} \rightarrow \text{Evidence} \rightarrow \text{Epistemic Verdict}
\end{align*}

We introduce a formal 7-tuple ontology $M_i = (R_i, O_i, C_i, P_i, E_i, A_i, V_i)$ decoupling statutory clauses from technical program properties. Crucially, Reg2App enforces an explicit three-tier conceptual distinction: (1) enacted parliamentary statutes and delegated administrative directives issued under statutory authority; (2) operationalized supervisory baseline controls derived from regulatory circulars and standards (e.g., OWASP MASVS); and (3) empirical security-risk indicators observed in client bytecode. We establish a two-dimensional observability taxonomy classifying requirements by level $\mathcal{L} \in \{\text{FULL}, \text{PARTIAL}, \text{NONE}\}$ and modality $\mathcal{M} \in \{\text{Static}, \text{Dynamic}, \text{Hybrid}, \text{Backend}, \text{Org}, \text{Legal}\}$. While this taxonomy conceptually encompasses all modalities, our implementation focuses on static Dalvik bytecode and manifest analysis. At the evaluation stage, Reg2App replaces scalar percentages with a discrete 5-state epistemic assessment space $\mathcal{S}$ preserving complete provenance. To eliminate linguistic barriers, Reg2App incorporates an evidence-grounded bilingual reporting engine generating actionable English and Bengali diagnostic summaries directly linked to statutory clauses.

\subsection{Research Questions}
\label{sec:intro:rqs}

This investigation addresses four research questions:
\textbf{RQ1 (Formalization and Observability):} How can statutory requirements be systematically formalized into machine-observable program properties, and how do client-side observational boundaries partition operationalized security controls?
\textbf{RQ2 (Detection Accuracy):} How accurately does Reg2App detect defined program properties on a controlled micro-benchmark suite, and how do its verdicts compare against independent manual analysis of production applications?
\textbf{RQ3 (Ecosystem Prevalence):} What is the empirical prevalence and sectoral distribution of regulatory non-conformance across an exploratory sample of production Android applications in Bangladesh?
\textbf{RQ4 (Usable Reporting):} Does evidence-grounded bilingual reporting improve developer comprehension and defect remediation efficiency compared to conventional scanner outputs?

\subsection{Key Contributions}
\label{sec:intro:contributions}

The primary contribution of this paper is a formally specified, provenance-preserving methodology for translating regulatory requirements into partially observable software properties. Specifically, we make three contributions:

\paragraph{1) Formal Methodology and Epistemic Observability Model} We establish an evidence-grounded mapping model $M_i$ with an explicit normative tier attribute $\tau_i$ and a two-dimensional observability operator $\Omega(R_i)$. We define a 5-state epistemic assessment space $\mathcal{S}$ with two objective metrics—Evidence Coverage ($EC$) and Evidence-Supported Compliance ($ESC$)—that strictly prevent compliance hallucinations on unobservable requirements. We evaluate 12 representative mappings across three statutory instruments through an expert-elicitation pilot of legal and security scholars (pre-reconciliation raw agreement $88.9\%$, $\text{Fleiss' } \kappa = 0.776$, $\text{Krippendorff's } \alpha = 0.782$, reaching 100\% unanimous consensus via structured Delphi reconciliation).

\paragraph{2) Static Analysis Engine and Provenance Architecture} We design and implement Reg2App, a static analysis pipeline inspecting Android Dalvik bytecode and manifest configurations. The engine constructs an Evidence Provenance Graph $G = (V, E)$, linking statutory articles to bytecode instructions, call graphs, and configuration snippets, ensuring unobservable requirements are deterministically partitioned.

\paragraph{3) Multi-Tiered Empirical Evaluation} We evaluate the methodology across four complementary dimensions: (i) an implementation correctness verification on 30 compilable micro-testcases verifying syntactic predicate extraction against synthetic code, accompanied by independent manual adjudication on production bytecode comprising a finding-sample precision study ($n=50$, $96.0\%$ precision [$48/50$, $95\%$ Wilson CI: $86.5\%$--$98.9\%$]) and an independent ground-truth control recall study ($n=51$, overall recall $94.1\%$ [$48/51$, $83.8\%$--$97.9\%$], defect recall $85.0\%$ [$17/20$]) across five applications; (ii) an empirical ecosystem study of 53 production Android applications across eight economic sectors sampled in early 2026, capturing an observational baseline disparity between heavily regulated financial applications and consumer retail platforms ($U = 156.0, p < 0.001$), shaped by unequal supervisory scopes and organizational compliance resources; (iii) an in-depth empirical audit of 39 production financial applications directly downloaded from Google Play, comprising all 14 operational Mobile Financial Services (MFS) applications and 25 commercial banking applications in Bangladesh, conducted under strict offline ethical auditing and institutional pseudonymization; and (iv) a counterbalanced human-subjects study ($N=60$) demonstrating that evidence-grounded bilingual reporting significantly improves developer comprehension and accelerates defect remediation ($52.1\%$ repair time reduction, from $39.7$ to $19.0$ minutes) over unparsed scanner traces. Bridging formal specification, static analysis, production audits, and human usability provides the end-to-end empirical evidence required to validate Regulation-as-Code systems in practice. All ontologies, benchmarks, and pipelines are released for open reproduction (\url{https://github.com/kh0kamoni/Reg2App}).

\subsection{Paper Organization}
\label{sec:intro:org}

Section~\ref{sec:motivating} presents a motivating example. Section~\ref{sec:statutory} details Bangladesh's statutory corpus. Section~\ref{sec:formal_framework} presents the formal 7-tuple mapping model and observability taxonomy (RQ1). Section~\ref{sec:expert} details the expert validation protocol. Section~\ref{sec:engine} describes the static analysis engine and provenance graph. Section~\ref{sec:assessment} formalizes the epistemic assessment space and Algorithm 1. Section~\ref{sec:benchmark} evaluates analyzer correctness on the micro-benchmark suite and ablation study (RQ2). Section~\ref{sec:ecosystem} presents the cross-sectoral ecosystem study (RQ3). Section~\ref{sec:case_study} details the empirical audit of production MFS and commercial banking applications. Section~\ref{sec:usable} reports the human-subjects study (RQ4). Section~\ref{sec:discussion} discusses implications and future work. Section~\ref{sec:threats} delineates threats to validity. Section~\ref{sec:related} reviews related work, and Section~\ref{sec:conclusion} concludes.
